\import{topology/topological-ring-power-bounded.tex}
\import{topology/topological-ring-uniformity.tex}
\import{topology/sequence-convergence.tex}

\section{Topologically nilpotent elements}

The power sequence is kept inline.  This avoids introducing a second public
power operation beside the existing ring power.

\begin{proposition}\label{topological_ring_power_sequence_type}
    Let $R$ be a topological ring and $x\in\carrier[R]$.
    Then
        $\Apply{\graphOnOp}{
            \naturals,\carrier[R],
            \fun{n:\Set}{\ringPower{R}{x}{n}}}
        \in\funs{\naturals}{\carrier[R]}$.
\end{proposition}
\begin{proof}
    $R$ is a ring by \obv{topological_ring_is_ring}.
    For every $n\in\naturals$ we have
        $\ringPower{R}{x}{n}\in\carrier[R]$ by \obv{ring_power_type}.
    Follows by \obv{logical_graph_on_function}.
\end{proof}

\begin{proposition}\label{topological_ring_power_sequence_apply}
    Let $R$ be a topological ring and $x\in\carrier[R]$.
    Let $n$ be a natural number.
    Then
        $(\Apply{\graphOnOp}{
            \naturals,\carrier[R],
            \fun{k:\Set}{\ringPower{R}{x}{k}}})(n)
        =\ringPower{R}{x}{n}$.
\end{proposition}
\begin{proof}
    Follows by \auto{logical_graph_on_apply}.
\end{proof}

\begin{definition}\label{topologically_nilpotent_element}
    $x$ is topologically nilpotent in $R$ iff
        $R$ is a topological ring and $x\in\carrier[R]$ and
        $\topologicalSequenceConverges{R}{
            \Apply{\graphOnOp}{
                \naturals,\carrier[R],
                \fun{n:\Set}{\ringPower{R}{x}{n}}}}{
            \ringZero[R]}$.
\end{definition}

\begin{proposition}\label{topologically_nilpotent_element_iff}
    Let $R$ be a topological ring and $x\in\carrier[R]$.
    Then $x$ is topologically nilpotent in $R$ iff
        $\topologicalSequenceConverges{R}{
            \Apply{\graphOnOp}{
                \naturals,\carrier[R],
                \fun{n:\Set}{\ringPower{R}{x}{n}}}}{
            \ringZero[R]}$.
\end{proposition}
\begin{proof}
    Follows by \auto{topologically_nilpotent_element}.
\end{proof}

\begin{proposition}\label{topologically_nilpotent_element_ambient}
    Suppose $x$ is topologically nilpotent in $R$.
    Then $R$ is a topological ring and $x\in\carrier[R]$.
\end{proposition}
\begin{proof}
    Follows by \obv{topologically_nilpotent_element}.
\end{proof}

\begin{theorem}\label{topologically_nilpotent_iff_additive_uniform_convergence}
    Let $R$ be a topological ring and $x\in\carrier[R]$.
    Then $x$ is topologically nilpotent in $R$ iff
        $\uniformSequenceConverges{\AdditiveUniformSpace{R}}{
            \Apply{\graphOnOp}{
                \naturals,\carrier[R],
                \fun{n:\Set}{\ringPower{R}{x}{n}}}}{
            \ringZero[R]}$.
\end{theorem}
\begin{proof}
    Let $U=\AdditiveUniformSpace{R}$.
    Let $s=\Apply{\graphOnOp}{
        \naturals,\carrier[R],
        \fun{n:\Set}{\ringPower{R}{x}{n}}}$.
    $U$ is a uniform space by
        \obv{additive_uniform_space_is_uniform_space}.
    Hence $\uniformTopology{U}$ is a topological space by
        \obv{uniform_topology_is_topological_space}.
    $\carrier[U]=\carrier[R]$ and
        $\carrier[\uniformTopology{U}]=\carrier[R]$ by
        \auto{additive_uniform_space_carrier,
            additive_uniform_topology_carrier}.
    $\opens[\uniformTopology{U}]=\opens[R]$ by
        \auto{additive_uniform_topology_opens}.
    $s\in\funs{\naturals}{\carrier[R]}$ by
        \auto{topological_ring_power_sequence_type}.
    $\ringZero[R]\in\carrier[R]$ by
        \obv{topological_ring_is_ring,ring_zero_type}.
    Hence $s\in\funs{\naturals}{\carrier[U]}$ and
        $\ringZero[R]\in\carrier[U]$.
    $\topologicalSequenceConverges{R}{s}{\ringZero[R]}$ iff
        $\topologicalSequenceConverges{\uniformTopology{U}}{s}{
            \ringZero[R]}$ by
        \obv{topological_sequence_convergence_invariant}.
    $\uniformSequenceConverges{U}{s}{\ringZero[R]}$ iff
        $\topologicalSequenceConverges{\uniformTopology{U}}{s}{
            \ringZero[R]}$ by
        \obv{uniform_sequence_convergence_iff_topological}.
    Follows by \auto{topologically_nilpotent_element}.
\end{proof}

\begin{theorem}\label{topologically_nilpotent_iff_eventually_in_neighbourhoods}
    Let $R$ be a topological ring and $x\in\carrier[R]$.
    Then $x$ is topologically nilpotent in $R$ iff
        for every $U\in\neighbourhoods{\ringZero[R]}{R}$ there exists
            $N\in\naturals$ such that for every $n\in\naturals$ such that
            $N\subseteq n$ we have $\ringPower{R}{x}{n}\in U$.
\end{theorem}
\begin{proof}
    Let $s=\Apply{\graphOnOp}{
        \naturals,\carrier[R],
        \fun{n:\Set}{\ringPower{R}{x}{n}}}$.
    $R$ is a topological space by
        \obv{topological_ring_is_topological_space}.
    $s\in\funs{\naturals}{\carrier[R]}$ by
        \auto{topological_ring_power_sequence_type}.
    $\ringZero[R]\in\carrier[R]$ by
        \obv{topological_ring_is_ring,ring_zero_type}.
    For every $n\in\naturals$ we have
        $s(n)=\ringPower{R}{x}{n}$ by
        \auto{topological_ring_power_sequence_apply}.
    Follows by
        \auto{topologically_nilpotent_element,
            topological_sequence_converges_iff_eventually_in_neighbourhoods}.
\end{proof}

\begin{theorem}\label{topologically_nilpotent_is_power_bounded}
    Let $R$ be a commutative topological ring.
    Suppose $x$ is topologically nilpotent in $R$.
    Then $x$ is power bounded in $R$.
\end{theorem}
\begin{proof}
    $R$ is a topological ring and $x\in\carrier[R]$ by
        \obv{topologically_nilpotent_element_ambient}.
    $R$ is a commutative ring by
        \obv{commutative_topological_ring_is_commutative_ring}.
    Hence $R$ is a ring by \obv{commutative_ring_is_ring}.
    Let $M=\MultiplicativeMonoid{R}$.
    Let $s=\Apply{\graphOnOp}{
        \naturals,\carrier[R],
        \fun{n:\Set}{\ringPower{R}{x}{n}}}$.
    Let $B=\pointwiseSubsetPowers{M}{\{x\}}$.
    $M$ is a monoid by \obv{multiplicative_monoid_is_exact_monoid}.
    $\carrier[M]=\carrier[R]$ and $\mul[M]=\ringMul[R]$ by
        \auto{multiplicative_monoid_fields}.
    Hence $x\in\carrier[M]$.
    $s\in\funs{\naturals}{\carrier[R]}$ by
        \auto{topological_ring_power_sequence_type}.
    Hence $s$ is a function and $\dom{s}=\naturals$ by
        \auto{funs_elim}.
    For every $n\in\naturals$ we have
        $s(n)=\ringPower{R}{x}{n}$ by
        \auto{topological_ring_power_sequence_apply}.
    $\{x\}\subseteq\carrier[M]$ by \auto{subseteq,singleton_iff}.
    Hence $B\subseteq\carrier[M]$ by
        \auto{pointwise_subset_powers_type}.
    Hence $B\subseteq\carrier[R]$.
    Show that for every $U\in\neighbourhoods{\ringZero[R]}{R}$ there
        exists $Q\in\neighbourhoods{\ringZero[R]}{R}$ such that
        $\pointwiseProduct{M}{Q}{B}\subseteq U$.
    \begin{subproof}
        Fix $U\in\neighbourhoods{\ringZero[R]}{R}$.
        Take $V,O\in\neighbourhoods{\ringZero[R]}{R}$ such that
            $\pointwiseProduct{\MultiplicativeMonoid{R}}{V}{O}
                \subseteq U$ by
            \obv{topological_ring_zero_neighbourhood_product}.
        Hence $\pointwiseProduct{M}{V}{O}\subseteq U$.
        Take $N\in\naturals$ such that for every $n\in\naturals$ such
            that $N\subseteq n$ we have $\ringPower{R}{x}{n}\in O$ by
            \auto{topologically_nilpotent_iff_eventually_in_neighbourhoods}.
        Let $F=\img{s}{N}$.
        $N$ is finite by \obv{natural_number_is_finite}.
        $N\subseteq\naturals$ by
            \auto{naturals_is_transitive_set,transitiveset_iff_subseteq}.
        Hence $N\subseteq\dom{s}$.
        Hence $\img{s}{N}$ is finite by \obv{finite_function_image}.
        Hence $F$ is finite.
        Show that every element of $F$ is an element of $\carrier[R]$.
        \begin{subproof}
            Fix $a\in F$.
            Hence $a\in\img{s}{N}$.
            Take $n\in N$ such that $(n,a)\in s$ by \obv{img_iff}.
            $n\in\dom{s}$ by \obv{subseteq}.
            Hence $a=s(n)$ by \auto{function_apply_iff}.
            $n\in\naturals$ by
                \obv{naturals_is_transitive_set,
                    transitiveset_iff_subseteq,subseteq}.
            Hence $s(n)\in\carrier[R]$ by \obv{funs_type_apply}.
        \end{subproof}
        Hence $F\subseteq\carrier[R]$ by \obv{subseteq}.
        Hence $F$ is left bounded in $R$ by
            \obv{finite_subset_left_bounded}.
        Take $W\in\neighbourhoods{\ringZero[R]}{R}$ such that
            $\pointwiseProduct{\MultiplicativeMonoid{R}}{W}{F}
                \subseteq U$ by
            \obv{left_bounded_in_neighbourhood}.
        Hence $\pointwiseProduct{M}{W}{F}\subseteq U$.
        Let $Q=V\inter W$.
        $V,W\subseteq\carrier[R]$ by
            \auto{topological_ring_zero_neighbourhoods_is_filter,filter}.
        Hence $Q\in\neighbourhoods{\ringZero[R]}{R}$ by
            \auto{topological_ring_zero_neighbourhoods_is_filter,
                filter_inter_in_iff}.
        $V\inter W\subseteq V$ and $V\inter W\subseteq W$ by
            \obv{inter,subseteq}.
        Hence $Q\subseteq V$ and $Q\subseteq W$.
        Show that every element of $\pointwiseProduct{M}{Q}{B}$ is an
            element of $U$.
        \begin{subproof}
            Fix $z\in\pointwiseProduct{M}{Q}{B}$.
            Take $q,b$ such that $q\in Q$ and $b\in B$ and
                $z=\mul[M](q,b)$ by \obv{pointwise_product_iff}.
            Hence $q\in V$ and $q\in W$ by \obv{subseteq}.
            Take a natural number $n$ such that
                $b\in\pointwiseSubsetPower{M}{\{x\}}{n}$ by
                \obv{pointwise_subset_powers_iff}.
            $\pointwiseSubsetPower{M}{\{x\}}{n}
                =\{\monoidPower{M}{x}{n}\}$ by
                \obv{pointwise_subset_power_singleton}.
            Hence $b=\monoidPower{M}{x}{n}$ by \auto{singleton_iff}.
            $\ringPower{R}{x}{n}=\monoidPower{M}{x}{n}$ by
                \obv{ring_power}.
            Hence $b=\ringPower{R}{x}{n}$.
            $s(n)=\ringPower{R}{x}{n}$.
            Hence $b=s(n)$.
            \begin{byCase}
                \caseOf{$N\subseteq n$.}
                    $\ringPower{R}{x}{n}\in O$.
                    Hence $b\in O$.
                    $\mul[M](q,b)\in\pointwiseProduct{M}{V}{O}$ by
                        \auto{pointwise_product_iff}.
                    Hence $\mul[M](q,b)\in U$ by \obv{subseteq}.
                \caseOf{$\lnot(N\subseteq n)$.}
                    $N$ and $n$ are ordinals by
                        \obv{natural_number_is_ordinal}.
                    $N\in n$ or $n\in N$ or $N=n$ by
                        \obv{ordinal_in_trichotomy}.
                    Hence $n\in N$ by
                        \obv{ordinal_elem_implies_subseteq,subseteq_refl}.
                    Hence $s(n)\in\img{s}{N}$ by
                        \auto{img_iff,function_apply_iff}.
                    Hence $b\in F$.
                    $\mul[M](q,b)\in\pointwiseProduct{M}{W}{F}$ by
                        \auto{pointwise_product_iff}.
                    Hence $\mul[M](q,b)\in U$ by \obv{subseteq}.
            \end{byCase}
        \end{subproof}
        Hence $\pointwiseProduct{M}{Q}{B}\subseteq U$ by \obv{subseteq}.
    \end{subproof}
    Hence $B$ is left bounded in $R$ by \auto{left_bounded_in}.
    Hence $\Bounded{R}{B}$ by
        \obv{bounded_in_commutative_topological_ring}.
    $\PowerBounded{R}{\{x\}}$ by \obv{power_bounded_set}.
    Follows by \obv{power_bounded_element}.
\end{proof}
